%% kapa: exam=1 round=33 year=2022 period=2 form=A subject=appraisal_law subject_ko=감정평가관계법규 seq=28 number=28
%% source: Q-Net 감정평가사 기출문제 (cst00309) artlSeq=5209997 file=raw/1008/5209997/2022년도 제33회 감정평가사 1차시험 2교시(감정평가 관계법규, 회계학).pdf
%% conversion: kordoc 4.24.51 -> tools/gen1.py (auto); review=proofread-sonnet; flags=none
\begin{question}{28}
공간정보의 구축 및 관리 등에 관한 법령상 토지와 지목이 옳게 연결된 것은?
\begin{choices}
  \item 묘지의 관리를 위한 건축물의 부지 - 묘지
  \item 원상회복을 조건으로 흙을 파는 곳으로 허가된 토지 - 잡종지
  \item 학교의 교사(校舍)와 이에 접속된 체육장 등 부속시설물의 부지 - 학교용지
  \item 자동차 판매 목적으로 설치된 야외전시장의 부지 - 주차장
  \item 자연의 유수가 있을 것으로 예상되는 소규모 수로부지 - 하천
\end{choices}
\end{question}
\answer{3}
